
\documentclass{article}
\usepackage{colortbl}
\usepackage{makecell}
\usepackage{multirow}
\usepackage{supertabular}

\begin{document}

\newcounter{utterance}

\centering \large Interaction Transcript for game `hot\_air\_balloon', experiment `air\_balloon\_survival\_en\_complexity\_easy', episode 1 with mistral{-}3{-}large{-}2512{-}t1.0.
\vspace{24pt}

{ \footnotesize  \setcounter{utterance}{1}
\setlength{\tabcolsep}{0pt}
\begin{supertabular}{c@{$\;$}|p{.15\linewidth}@{}p{.15\linewidth}p{.15\linewidth}p{.15\linewidth}p{.15\linewidth}p{.15\linewidth}}
    \# & $\;$A & \multicolumn{4}{c}{Game Master} & $\;\:$B\\
    \hline

    \theutterance \stepcounter{utterance}  
    & & \multicolumn{4}{p{0.6\linewidth}}{
        \cellcolor[rgb]{0.9,0.9,0.9}{
            \makecell[{{p{\linewidth}}}]{
                \texttt{\tiny{[P1$\langle$GM]}}
                \texttt{You are participating in a collaborative negotiation game.} \\
\\ 
\texttt{Together with another participant, you must agree on a single set of items that will be kept. Each of you has your own view of how much each item matters to you (importance). You do not know how the other participant values the items. Additionally, you are given the effort each item demands.} \\
\texttt{You may only agree on a set if the total effort of the selected items does not exceed a shared limit:} \\
\\ 
\texttt{LIMIT = 8170} \\
\\ 
\texttt{Here are the individual item effort values:} \\
\\ 
\texttt{Item effort = \{"B21": 594, "B00": 551, "A08": 634, "A27": 204, "B94": 557, "B53": 440, "A23": 734, "C65": 893, "A89": 246, "B67": 592, "B45": 143, "C34": 568, "C12": 472, "A05": 401, "B60": 729, "C92": 201, "B69": 85, "A84": 641, "C06": 958, "B12": 80, "A16": 157, "B85": 806, "C35": 684, "B10": 59, "C15": 31, "B09": 765, "B16": 416, "C81": 392, "A58": 428, "C24": 699, "A30": 141, "A66": 606, "B75": 612, "B46": 133, "C70": 689\}} \\
\\ 
\texttt{Here is your personal view on the importance of each item:} \\
\\ 
\texttt{Item importance values = \{"B21": 138, "B00": 583, "A08": 868, "A27": 822, "B94": 783, "B53": 65, "A23": 262, "C65": 121, "A89": 508, "B67": 780, "B45": 461, "C34": 484, "C12": 668, "A05": 389, "B60": 808, "C92": 215, "B69": 97, "A84": 500, "C06": 30, "B12": 915, "A16": 856, "B85": 400, "C35": 444, "B10": 623, "C15": 781, "B09": 786, "B16": 3, "C81": 713, "A58": 457, "C24": 273, "A30": 739, "A66": 822, "B75": 235, "B46": 606, "C70": 105\}} \\
\\ 
\texttt{Goal:} \\
\\ 
\texttt{Your goal is to negotiate a shared set of items that benefits you as much as possible (i.e., maximizes total importance to YOU), while staying within the LIMIT. You are not required to make a PROPOSAL in every message {-} you can simply negotiate as well. All tactics are allowed!} \\
\\ 
\texttt{Interaction Protocol:} \\
\\ 
\texttt{You may only use the following structured formats in a message:} \\
\\ 
\texttt{PROPOSAL: \{'A', 'B', 'C', …\}} \\
\texttt{Propose keeping exactly those items.} \\
\texttt{REFUSE: \{'A', 'B', 'C', …\}} \\
\texttt{Explicitly reject opponent's proposal.} \\
\texttt{ARGUMENT: \{'...'\}} \\
\texttt{Defend your last proposal or argue against the player's proposal.} \\
\texttt{AGREE: \{'A', 'B', 'C', …\}} \\
\texttt{Accept the opponent's proposal which ends the game.} \\
\texttt{STRATEGIC REASONING: \{'...'\}} \\
\texttt{	Describe your strategic reasoning or anticipation explaining your choice of action. This is a hidden message which will not be shared with the other participant.} \\
\\ 
\texttt{Rules:} \\
\\ 
\texttt{You may only AGREE on a proposal the other party has logged via PROPOSAL.} \\
\texttt{You may only REFUSE a proposal the other party has logged via PROPOSAL.} \\
\texttt{Total effort of any PROPOSAL or AGREE set must be ≤ LIMIT.} \\
\texttt{Do NOT reveal your hidden importance scores.} \\
\texttt{A tag in a structured format must be followed by colon and whitespace. The argument must be a python set containing 0 or more strings.} \\
\texttt{So, it must be of the form TAG: \{...\}} \\
\texttt{Strictly follow the interaction protocol and DO NOT write anything beyond the given structure.} \\
\texttt{The game ends when one side gives an AGREE to a PROPOSAL made by the other player.} \\
\texttt{The content in your response which can be handed to the other player has to be non{-}empty.} \\
\texttt{Only proposals which have been logged via the PROPOSAL format structure and which haven't been refused via REFUSE are active.} \\
\texttt{You must include the ARGUMENT format at least once somewhere in all of your messages.} \\
\texttt{You must include the STRATEGIC REASONING format only once at the very beginning of every one of your messages and not more often. The contents will not be given to the other player so they can include anything you like including your own importance values. Here you should reason step by step to come up with you next move.} \\
            }
        }
    }
    & \\ \\

    \theutterance \stepcounter{utterance}  
    & \multicolumn{4}{p{0.6\linewidth}}{
        \cellcolor[rgb]{0.9,0.9,0.9}{
            \makecell[{{p{\linewidth}}}]{
                \texttt{\tiny{[P1$\rangle$GM]}}
                \texttt{STRATEGIC REASONING: \{} \\
\texttt{To maximize my total importance while staying under the LIMIT of 8170, I need to prioritize items with high importance to me and low effort. Here's my step{-}by{-}step reasoning:} \\
\\ 
\texttt{1. **Initial Analysis**:} \\
\texttt{   {-} Sort items by importance{-}to{-}effort ratio (importance/effort) to identify the most "efficient" items.} \\
\texttt{   {-} High{-}importance, low{-}effort items are ideal (e.g., "B12": 915 importance, 80 effort; "C15": 781 importance, 31 effort).} \\
\texttt{   {-} Avoid items with low importance and high effort (e.g., "C06": 30 importance, 958 effort; "B53": 65 importance, 440 effort).} \\
\\ 
\texttt{2. **Greedy Selection**:} \\
\texttt{   {-} Start with the highest importance{-}to{-}effort items and add them until the LIMIT is reached.} \\
\texttt{   {-} Example greedy set (not necessarily optimal but a strong starting point):} \\
\texttt{     {-} Top items: "B12", "C15", "B10", "B46", "A16", "A30", "B45", "C81", "A27", "A08", "B00", "B60", "B94", "B67", "C12", "A89", "A66".} \\
\texttt{     {-} Total effort: \textasciitilde{}7900 (under LIMIT), total importance: \textasciitilde{}10,000+.} \\
\texttt{   {-} This is likely too aggressive; the opponent may refuse. Need to negotiate.} \\
\\ 
\texttt{3. **Negotiation Strategy**:} \\
\texttt{   {-} Start with a proposal that includes many high{-}importance items but leaves room for compromise.} \\
\texttt{   {-} Use ARGUMENT to justify why certain items are critical (without revealing exact importance).} \\
\texttt{   {-} If the opponent refuses, adjust by swapping out items with lower importance{-}to{-}effort ratios.} \\
\texttt{   {-} Avoid revealing my exact priorities; focus on "mutual benefit" language.} \\
\\ 
\texttt{4. **First Proposal**:} \\
\texttt{   {-} Propose a set with a mix of high{-}importance (to me) and moderate{-}effort items.} \\
\texttt{   {-} Include some items that might be valuable to the opponent (e.g., "A08", "A27", "B60") to} \\
            }
        }
    }
    & & \\ \\

    \theutterance \stepcounter{utterance}  
    & & \multicolumn{4}{p{0.6\linewidth}}{
        \cellcolor[rgb]{0.9,0.9,0.9}{
            \makecell[{{p{\linewidth}}}]{
                \texttt{\tiny{[P1$\langle$GM]}}
                \texttt{Your response did not start with the proper strategic reasoning tag at the very beginning of your response. The very first structured format must be of the form STRATEGIC REASONING: \{...\}. Try again.} \\
            }
        }
    }
    & \\ \\

    \theutterance \stepcounter{utterance}  
    & \multicolumn{4}{p{0.6\linewidth}}{
        \cellcolor[rgb]{0.9,0.9,0.9}{
            \makecell[{{p{\linewidth}}}]{
                \texttt{\tiny{[P1$\rangle$GM]}}
                \texttt{STRATEGIC REASONING: \{} \\
\texttt{1. **Objective**: Maximize my total importance while keeping total effort ≤ 8170. I must negotiate a set that includes as many high{-}importance (to me) items as possible, ideally without revealing my exact priorities.} \\
\texttt{2. **Initial Prioritization**:} \\
\texttt{   {-} Sort items by importance{-}to{-}effort ratio (descending) to identify the most "valuable" items for me.} \\
\texttt{   {-} Top candidates (high importance, low effort):} \\
\texttt{     {-} "B12" (915/80 = 11.44), "C15" (781/31 = 25.19), "B10" (623/59 = 10.56), "B46" (606/133 = 4.56), "A16" (856/157 = 5.45), "A30" (739/141 = 5.24), "B45" (461/143 = 3.22), "C81" (713/392 = 1.82).} \\
\texttt{   {-} Also include high{-}importance items even if effort is moderate (e.g., "A08" (868/634), "B00" (583/551), "B60" (808/729)).} \\
\texttt{3. **First Proposal**:} \\
\texttt{   {-} Start with a set that includes all top candidates and leaves \textasciitilde{}1000 effort buffer for negotiation.} \\
\texttt{   {-} Proposed set (effort: 7170, importance: 8500+):} \\
\texttt{     \{'B12', 'C15', 'B10', 'B46', 'A16', 'A30', 'B45', 'C81', 'A27', 'A08', 'B00', 'B60', 'B94', 'B67', 'C12'\}.} \\
\texttt{   {-} This is aggressive but leaves room to swap out items if the opponent refuses.} \\
\texttt{4. **Negotiation Tactics**:} \\
\texttt{   {-} Use ARGUMENT to emphasize "mutual benefit" (e.g., "These items} \\
            }
        }
    }
    & & \\ \\

    \theutterance \stepcounter{utterance}  
    & & \multicolumn{4}{p{0.6\linewidth}}{
        \cellcolor[rgb]{0.9,0.9,0.9}{
            \makecell[{{p{\linewidth}}}]{
                \texttt{\tiny{[P1$\langle$GM]}}
                \texttt{Your response did not start with the proper strategic reasoning tag at the very beginning of your response. The very first structured format must be of the form STRATEGIC REASONING: \{...\}. Try again.} \\
            }
        }
    }
    & \\ \\

    \theutterance \stepcounter{utterance}  
    & \multicolumn{4}{p{0.6\linewidth}}{
        \cellcolor[rgb]{0.9,0.9,0.9}{
            \makecell[{{p{\linewidth}}}]{
                \texttt{\tiny{[P1$\rangle$GM]}}
                \texttt{STRATEGIC REASONING: \{} \\
\texttt{1. **Re{-}evaluating Approach**: My initial proposal was too aggressive and may be rejected outright. I need a more balanced first offer that still prioritizes my high{-}importance items while appearing reasonable to the opponent.} \\
\texttt{2. **Revised Prioritization**:} \\
\texttt{   {-} Focus on items with the highest importance{-}to{-}effort ratio first (e.g., "B12", "C15", "B10", "B46", "A16").} \\
\texttt{   {-} Include some high{-}importance items even if their ratio is lower (e.g., "A08", "B60", "B94") to anchor my preferences.} \\
\texttt{   {-} Exclude low{-}importance, high{-}effort items (e.g., "C06", "B53", "C65") to save effort.} \\
\texttt{3. **First Proposal**:} \\
\texttt{   {-} Propose a set with total effort \textasciitilde{}7000 (leaving \textasciitilde{}1170 buffer for negotiation).} \\
\texttt{   {-} Items: \{'B12', 'C15', 'B10', 'B46', 'A16', 'A30', 'B45', 'A27', 'A08', 'B00', 'B60', 'B67', 'C12'\}.} \\
\texttt{   {-} Total effort: 59 + 31 + 59 + 133 + 157 + 141 + 143 + 204 + 634 + 551 + 729 + 592 + 472 = 3905 (too low; needs adjustment).} \\
\texttt{   {-} Adjust to include more high{-}importance items: Add 'B94' (557) and 'A89' (246).} \\
\texttt{   {-} New total effort: 3905 + 557 + 246 = 4708 (still too low).} \\
\texttt{   {-} Final adjustment: Add 'C81' (392) and 'A66' (606).} \\
\texttt{   {-} Final total effort: 4708 + 392 + 606 = 5706 (importance: \textasciitilde{}8000).} \\
\texttt{   {-} This is still too conservative; aim} \\
            }
        }
    }
    & & \\ \\

    \theutterance \stepcounter{utterance}  
    & & & \multicolumn{2}{p{0.3\linewidth}}{
        \cellcolor[rgb]{0.9,0.9,0.9}{
            \makecell[{{p{\linewidth}}}]{
                \texttt{\tiny{[GM$|$GM]}}
                \texttt{abort} \\
            }
        }
    }
    & & \\ \\

    \theutterance \stepcounter{utterance}  
    & & & \multicolumn{2}{p{0.3\linewidth}}{
        \cellcolor[rgb]{0.9,0.9,0.9}{
            \makecell[{{p{\linewidth}}}]{
                \texttt{\tiny{[GM$|$GM]}}
                \texttt{end game} \\
            }
        }
    }
    & & \\ \\

\end{supertabular}
}

\end{document}
